\begin{description}
\item[(a)] Take $\lambda_S=180^\circ$ on 22 September (also accept 23 September). With $D$ the number of days since the chosen equinox,
\[
\lambda_S=0^\circ/180^\circ+\frac{D}{365.25}\times360^\circ.
\]\hfill(0.5 point)
The two choices give $\lambda_S=207.6^\circ$ or $206.6^\circ$. \hfill(0.5 point)
\[
\sin\delta=\sin\beta\cos\varepsilon+\cos\beta\sin\varepsilon\cos(90^\circ-\lambda)
=\sin\varepsilon\sin\lambda\quad(\beta=0^\circ).
\]\hfill(1 point)
Therefore $\delta=-10^\circ38'$ or $-10^\circ16'$. \hfill(0.5 point)
Using $\cos\alpha=\cos\lambda/\cos\delta$ \hfill(0.5 point)
gives $\alpha=13^{\mathrm h}42^{\mathrm m}$ or $13^{\mathrm h}38^{\mathrm m}$. \hfill(1 point)
Award 75\% for computing $\alpha$ directly with uniform motion on the equator, and 25\% for Euclidean geometry in place of spherical trigonometry.
\item[(b)]
\begin{align*}
\cos(90^\circ-h)&=\cos(90^\circ-\varphi)\cos(90^\circ-\delta)\\
&\quad+\sin(90^\circ-\varphi)\sin(90^\circ-\delta)\cos H,\\
0&=\sin\varphi\sin\delta+\cos\varphi\cos\delta\cos H,\\
\cos H&=-\tan\varphi\tan\delta.
\end{align*}\hfill(2 points)
Writing the last formula directly receives full credit. Local time is
\[
t_l=H+12^{\mathrm h}+n-L+\eta+\mathrm{D.S.T.}
\]
The other terms are constant during a day, so $\Delta t_l\approx\Delta H$. \hfill(0.5 point)
The two solutions of $\cos x=a$ in $[0,2\pi]$ are $\arccos a$ and $2\pi-\arccos a$, giving $t=2\arccos(-\tan\varphi\tan\delta)$ after conversion to hours. The four durations are
\[
t_1=10^{\mathrm h}42^{\mathrm m},\quad t_2=12^{\mathrm h}35^{\mathrm m},\quad
 t_3=11^{\mathrm h}31^{\mathrm m},\quad t_4=10^{\mathrm h}16^{\mathrm m}.
\]\hfill(1 point)
Theodor is in city 1, Kutaisi. \hfill(0.5 point)
Full credit is also awarded for calculating both solutions for $H$ for each city and taking their time difference.
\item[(c)] For the Sun to be directly overhead, $h_S=90^\circ$ requires $\varphi=\delta$. \hfill(0.5 point)
Hence $\varphi=-10^\circ38'$. \hfill(0.5 point)
With $L=12^{\mathrm h}-T_{\mathrm{GMT}}$ \hfill(0.5 point)
and $T_{\mathrm{GMT}}=H+12^{\mathrm h}-L$, the longitude is $37^\circ47'$ W, or $37^\circ1'$ W using 23 September. \hfill(0.5 point)
\end{description}
